\section{서론}

Representation learning은 내부 상태를 vector나 tensor로 저장한다. 이러한 저장
방식만으로는 어떤 entity가 지속되는지, 어떤 relation이 합성되는지, 어떤 metric
quantity가 중요한지, 순서 및 질량 관측량이 어떻게 작동하는지, action이 상태를
어떻게 바꾸는지가 정해지지 않는다. TSI는 더 제한된 질문에서 출발한다. 내부
상태가 structural simulation, reasoning, action을 지원하려면 어떤 수학적 데이터가
명시되어야 하는가?

이 논문에서 전개하는 형식적 대상은 finite presented schema, common tagged
carrier, typed relation, simplicial complex, metric, filtration, full-support
mass, action-conditioned tracked transition으로 구성된다. 이 layer들은 상호작용할
수 있지만 서로 동일시되지는 않는다. Vector가 이 데이터를 encode할 수는 있으나,
그 좌표만으로 전체 structural specification이 정해지지는 않는다.

이 논문은 다섯 구성요소를 하나의 통합된 기여로 제시한다. (1) 선언된 cross-layer
compatibility axiom을 갖는 state contract, (2) exact structural isomorphism과
그 동치관계 증명, (3) 다섯 observable에 대한 survivor-level preservation과 그
합성 증명, (4) local preservation과 turnover를 분리하는 six-defect zero
criterion, (5) metric--measure comparison과 rollout stability를 포함한 category
및 quotient-level 확장이다.

Finite-relation category와 presented-category descent는 relation calculus 및
categorical data formalism에서 도입한다
\cite{freydscedrov1990categories,spivak2012functorial,
spivakwisnesky2015relational}. Filtered persistence stability는 persistent
homology에서 도입한다 \cite{cohensteiner2007stability,chazal2009proximity}.
Metric--measure coupling과 gluing은 Gromov--Wasserstein 및 optimal-transport
이론에서 도입한다 \cite{memoli2011gromov,peyre2019computational}. Rollout
recurrence는 표준 Lipschitz error-propagation pattern을 사용한다
\cite{asadi2018lipschitz}. TSI 고유의 기여는 이들을 common tagged state
contract 아래 통합하고, compatibility invariant와 label-preserving partial
transition category를 정의하며, survivor preservation과 temporal identity를
명시적으로 분리한 데 있다. 이 논문은 이러한 layer가 intelligence에 필요충분하다고,
human mathematical imagery를 기술한다고, 또는 neural model이 inductive
constraint 없이 이를 발견한다고 주장하지 않는다.
